\section{Running the System}
\label{sec:running}

\begin{table*}[t]
\centering
\scriptsize
\begin{verbatim}
# 0. Environment: Python 3.11, one CUDA GPU (CPU formats need ctranslate2 and whisper.cpp)
micromamba create -n cs425 -c conda-forge python=3.11 pip ffmpeg && micromamba activate cs425
pip install -r requirements.txt && pip install -e . && pytest        # normaliser + scorer tests
# 1. Pre-processing: 16 kHz mono FLAC clips + JSONL manifests in data/<corpus>/manifests/
huggingface-cli download CAiRE/ASCEND --repo-type dataset --local-dir data/ascend/hf
python scripts/prepare_ascend.py                       # official splits
python scripts/fetch_nsc_dropbox.py fetch data/nsc_part4 -j 2 --only cs-chn && python scripts/prepare_nsc.py
python scripts/prepare_fleurs.py                       # spliced clips for the language-ID head
python scripts/prepare_stress.py                       # our recordings: data/stress/audio/<reader>/<id>.wav
python scripts/align.py ascend_train nsc_train ascend_val nsc_val    # MMS word times
# 2. Training (checkpoints in ckpt/<run>/best)
python scripts/train.py configs/ft_base.yaml           # plain fine-tuning; also ft_tiny.yaml, ft_small.yaml
python scripts/train.py configs/ft_base.yaml --run full_base --alpha 0.5 --beta 0.5 --lambda_sem 0.1 --rho0 0.5
python scripts/train_lid.py --encoder ckpt/ft_base/best --run lid_ft_base
python scripts/train.py configs/ft_base.yaml --run qat_full_base --model ckpt/full_base/best \
       --lora true --lr 1e-4 --epochs 1 --warmup 100
# 3. Testing (predictions in preds/<model>__<precision>__<method>.jsonl)
python scripts/decode.py --model ckpt/full_base/best --model_tag full_base --precision nf4_bnb --method full
python scripts/deploy_eval.py --model ckpt/full_base/best --tag full_base --method full \
       --formats ct2_f32 ct2_int8 ggml_f16 ggml_q8_0 ggml_q4_0
python scripts/sensitivity.py --model ckpt/full_base/best --tag full_base      # mixed-precision blocks
python scripts/rerank.py --model ckpt/full_base/best --tag full_base --lid ckpt/lid_ft_base/head.pt --loop_filter
# 4. Scoring and tables (results/tables/*.csv, results/figures/)
python scripts/score.py full_base__fp16__full full_base__nf4_bnb__full     # all metrics, 95% CIs
python scripts/score.py --fragility full_base__fp16__full full_base__nf4_bnb__full --set nsc_test --main en
python scripts/make_results.py && python scripts/export_tables.py && python scripts/make_figures.py
GPU=0 bash scripts/queue_h.sh && python scripts/make_human_tables.py       # human stress recordings
\end{verbatim}
\vspace{-1em}
\caption{Commands to pre-process the data, train, test and score (run from the repository root; GPU chosen with \texttt{CUDA\_VISIBLE\_DEVICES}). The \texttt{scripts/queue\_*.sh} files chain these steps exactly as we ran them.}
\label{tab:commands}
\end{table*}

All code is in one repository (\url{https://github.com/nibess1/cs425_codeswitch}); Table~\ref{tab:commands} lists the commands. Each step reads and writes plain files, so steps can be rerun separately. \emph{Pre-processing} turns each corpus into 16\,kHz FLAC clips and a JSONL manifest with the raw and normalised reference, the speaker, the main language and the duration. \emph{Training} reads a YAML config whose fields can be overridden on the command line (Table~\ref{tab:hparams}). \emph{Testing} decodes the test sets with a 16-bit, simulated 8/4-bit (bitsandbytes) or exported CPU model (CTranslate2, whisper.cpp) and writes one JSONL line per utterance. \emph{Scoring} (\texttt{score.py}) computes every metric of \S\ref{sec:task} from those files with bootstrap intervals; every number in this report comes from it.
